\documentclass{article}
\usepackage{graphicx} % Required for inserting images
\usepackage{amssymb, amsmath}

\title{Filosofia da Matemática 
\\(F):'What is a categorical theory? What problem does the existence of non-categorical mathematical theories pose for Platonism in mathematics?'}
\author{106618 }
\date{Número de palavras: 1489}

\begin{document}
\newtheorem{dfn}{Definição}
\renewcommand{\refname}{Bibliografia}

\maketitle
A natureza dos objetos matemáticos e a forma como os termos matemáticos se referem a esses objetos constituem um dos problemas centrais da filosofia da matemática. Mesmo admitindo a existência de tais objetos matemáticos, permanece por esclarecer de que modo as nossas teorias e termos matemáticos conseguem referir precisamente a essas entidades.


No presente ensaio são consideradas duas variantes do platonismo e algumas das suas consequências. Começamos por explorar os platonistas modelistas e o seu problema com teorias não categóricas. De seguida, estudamos os platonistas moderados e a inevitável indeterminação da referência dos objetos matemáticos. Finalmente, apresentamos algumas propostas que tentam garantir a determinação da referência em matemática. \footnote{Nas instruções do ensaio, este tema era o único sem "Argue in favor of your position". Deste modo, prioritizou-se a exposição das várias teses, em detrimento da defesa de alguma delas.}

\vspace{1cm}

Começamos com algumas definições fundamentais de teoria dos modelos. \footnote{Em anexo segue a formalização da lógica de 1º ordem que é usada neste ensaio.}

\vspace{0.5cm}

\begin{dfn}
    Dada uma teoria $\Gamma$, um modelo de $\Gamma$ é uma estrutura de interpretação $I=(\mathcal{D}, \_^{\mathcal{F}}, \_^{\mathcal{P}})$ tal que $I \Vdash A$, para todo o $A \in \Gamma$.
\end{dfn}

\vspace{0.5cm}

\begin{dfn}
    Sejam $I_1=(\mathcal{D}_1,\_^{\mathcal{F}}_1,\_^{\mathcal{P}}_1), I_2=(\mathcal{D}_2,\_^{\mathcal{F}}_2,\_^{\mathcal{P}}_2)$ dois modelos de uma certa teoria. Um isomorfismo entre $I_1, I_2$ é uma função $h \in \mathcal{D}_1 \rightarrow\mathcal{D}_2$ tal que
    \begin{itemize}
        \item $h(c^{I{_1}})=c^{I_{2}}$, para qualquer $c \in \mathcal{F}_0$
        \item $h(f^{I_1}(t_1,...,t_n))=f^{I_2}(h(t_1),...,h(t_n))$, para qualquer $f \in \mathcal{F}_n$
        \item $(t_1,...,t_n) \in P^{I_1}$ sse $(h(t_1),...,h(t_n)) \in P^{I_2}$
    \end{itemize}
    Nesse caso, dizemos que $I\;e\;J$ são isomorfos e escrevemos $I \cong J$.
\end{dfn}

\vspace{0.5cm}

\begin{dfn}
    Uma teoria diz-se categórica se e só se todos os seus modelos forem isomórficos.
\end{dfn}

\vspace{0.5cm}
Intuitivamente, uma teoria categórica determina completamente a estrutura do seu objeto de estudo, uma vez que todos os modelos que a satisfazem são isomorfos entre si.

\vspace{1cm}

Para o platonista, os objetos matemáticos existem realmente, de forma objetiva. São entidades abstratas, residentes fora do espaço tempo, cuja existência é independente de mentes humanas.

No presente caso, estamos interessados numa variação desta corrente filosófica: o platonismo modelista, ou \textit{objects modelists} como em \cite{buttonwalsh2018} são referidos. Para além da tendência platonista, estes pensadores não se interessam pelos modelos individuais de uma teoria. Em vez disso, o seu interesse centra-se nas suas classes de equivalências, ou \textit{isomorphism types}. Este último conceito é definido considerando a seguinte relação. Dados dois modelos $I,J$ de $\Gamma$:

\begin{equation}
    I \sim J \;\;\; \text{sse}\;\;\; I \cong J
\end{equation}

É imediato verificar que esta relação é, de facto, uma relação de equivalência. Consequentemente, os tipos de isomorfismos são definidos como as classes de equivalência da relação anterior.

\vspace{0.5cm}

Observe-se ainda que uma teoria é categórica se, e só se, o conjunto dos seus tipos de isomorfismo contém, no máximo, um elemento.

\vspace{0.5cm}

Acontece que, para o platonista moderado, a propriedade de categoricidade de uma teoria possui implicações particularmente fortes.

Suponha-se que uma teoria $\Gamma$ não é categórica. Nesse caso, existem pelo menos dois modelos $N,M$ não isomorfos ($N \ncong M$). Consequentemente, $N$ e $M$ pertencem a classes de equivalência diferentes, denotemo-las por $\mathcal{M}, \mathcal{N}$, respetivamente. Surge então a questão natural, referida em \cite{buttonwalsh2018} como \textit{ Modelist’s Doxological Challenge}:
\begin{center}
    Qual dos tipos de isomorfismo $\mathcal{M}$ ou $\mathcal{N}$, deve ser considerado o correto? \footnote{Devido às limitações do ensaio, não é possível explorar algumas das respostas a esta questão. Refere-se o leitor para \cite{buttonwalsh2018}}
\end{center}

\vspace{2cm}

Concetremo-nos agora nos platonistas moderados. A componente platonista corresponde às ideias anteriormente apresentadas. Já a componente moderada está relacionada com o naturalismo no sentido em que rejeita uma faculdade especial de ”intuicão matemática”. Esta visão defende que a referência a objetos matemáticos não se faz "por contacto", mas sim "por descrição".

\vspace{0.5cm}

Note-se que, ao contrário, do platonista modelista, o platonista moderado interessa-se pelos modelos individuais de uma teoria e não pelos seus tipos de isomorfismos. Isto quer dizer que, se tivermos dois modelos distintos mas isomorfos, então existe pelo menos um objeto com referência diferente nos dois modelos, dando origem a uma situação de indeterminação da referência.

O passo seguinte consiste em compreender quando e como é possível construir dois modelos diferentes mas isomorfos de uma mesma teoria. O argumento modelo-teorético de Putnam produz exatamente isso. A ideia principal é que uma permutaçaõ, não trivial, do domínio do modelo consiste num isomorfismo do modelo original. Assim, assumindo que a teoria possui um modelo cujo domínio contém mais do que um objeto, é possível, permutando os elementos do domínio, obter um novo modelo distinto mas isomorfo ao original. Mais precisamente, o argumento pode ser apresentado como:

\begin{enumerate}
    \item platonismo moderado                                                                               \hfill (suposição)
    \item \parbox[t]{0.9\linewidth}{%
    os objetos referidos pelos termos de uma teoria são aqueles que satisfazem
    os axiomas da mesma%
    }\hfill(de 1)
    
    \item \parbox[t]{0.9\linewidth}{%
    qualquer permutação de um qualquer objeto satisfaz exatamente as mesmas verdades que o objeto original%
    }\hfill(Teo.)
    
    \item \parbox[t]{0.9\linewidth}{%
    existem permutações não triviais para todo o domínio de \\objetos, com cardinalidade maior que 1%
    }\hfill(Lógica)

    
    $\therefore$ Indeterminação radical da referência
\end{enumerate}

\vspace{0.5cm}

A referência é indeterminada no seguinte sentido. Consideremos uma teoria qualquer da aritmética. Segundo o argumento anterior, para todo o objeto do domínio existe um modelo, isomorfo ao modelo original, no qual esse objeto é a interpretação do termo $'1'$. Isto é, todo o termo refere-se igualmente a todo o objeto. Mais geralmente, o mesmo raciocínio aplica-se aos restantes termos matemáticos.

\vspace{0.5cm}

À primeira vista, pode parecer que a segunda premissa, decorrente da suposição do platonismo moderado, não foi utilizada na derivação da conclusão e que, por conseguinte, a tese da indeterminação radical da referência poderia ser generalizada a qualquer corrente da filosofia da matemática. 

Vejamos que esse não é o caso. Para tal, retiremos a suposição do platonismo moderado e apliquemos o argumento da permutação de Putnam. Uma resposta, ou objeção, clássica à indeterminação resultante do argumento anterior consiste em sustentar que os diferentes modelos obtidos possuem graus de preferência diferentes. Em particular, defende-se que existe um modelo privilegiado, ou o modelo "correto". Deste modo, a tese da indeterminação perderia força e o argumento deixaria de ser convincente. É neste ponto que entra em cena a segunda premissa. Para o platonista moderado, mais especificamente devido à componente moderada, a ideia de um modelo preferível é incompatível com os seus valores teóricos. Com efeito, esta visão rejeita a noção de uma intuição matemática capaz de identificar um modelo correto. É precisamente esta ideia que se encontra expressa na segunda permissa, o que explica a sua necessidade no argumento. 

Concluímos, assim, que o platonista moderado está comprometido com a tese da indeterminação radical da referência. Em contrapartida, um platonista não moderado, é capaz de defender a determinação da referência, recorrendo à ideia de um modelo privilegiado. O platonista confronta-se, assim, com o dilema de escolher entre a moderação e a determinação da referência.

\vspace{0.5cm}

Consideremos agora algumas tentativas de garantir a determinação da referência matemática.\\

O próprio Putnam responde, ainda que brevemente, ao seu argumento, abandonando o platonismo moderado em favor do chamado realismo interno. Trata-se de uma visão construtiva, no sentido em que asserta que os objetos matemáticos são construidos, o que a torna incompatível com o platonismo moderado. É precisamente durante esse processo construção que os objetos recebem as suas designações, fixando-se assim a referência dos termos matemáticos. 

Note-se que este é outra instância da ideia do modelo preferível. Neste caso, o modelo preferível é aquele cujas interpretações dos termos matemáticos coincidem com as designações estabelecidas pela construção dos mesmos.

\vspace{0.5cm}

Outra resposta é apresentada por Michael Devitt em \cite{devitt1984}. Para a compreender, importa recordar a resposta de Putnam à objeção dos modelos preferíveis, considerada sem a suposição do platonismo moderado. Segundo Putnam, tal objeção falha porque pressupõe uma teoria capaz de explicar qual seria o tal modelo privilegiado. No entanto, essa explicação constitui apenas "apenas mais teoria" e, portanto, os seus próprios termos estarão sujeitos ao mesmo problema de indeterminação. 

Devitt , enquanto realista, defende que os termos de uma teoria referem-se a objetos reais através de relações causais. Putnam responde perguntando o que fixa a referência de "causalmente relacionado". Para responder a esta questão, o realista formula uma nova teoria. Novamente, Putnam questiona sobre o que fixa a referência dos termos dessa nova teoria? E assim sucessivamente. Putnam defende que somos incapazes de chegar a uma relação que fixe a referência.

É neste instante que Devitt identifica aquilo que considera ser o erro de Putnam. Recorrendo a um exemplo ilustrativo. Uma criança pergunta: "Porque chove?". Uma explicação é-lhe dada. A criança insiste: "E porque tal acontece?". Fornece-se uma outra explicação, mas criança não se conforma: "E porque é esse o caso?". Este processo pode prosseguir indefinidamente. O importante é que o facto de a criança poder continuar a formular perguntas não demonstra que as explicações anteriores sejam falsas. Segundo Devitt, é precisamente este o erro que Putnam está a cometer em relação à referência.


\vspace{2cm}

Conclui-se, assim, que diferentes variantes do platonismo enfrentam diferentes problemas. Por um lado, o platonista modelista depara-se com dificuldades com teorias não categóricas. Por outro lado, o platonista moderado parece estar comprometido com a indeterminação radical da referência dos termos matemáticos, apenas evitada abandonando a componente moderada do pensamento. Finalmente, Putnam e Devitt demonstram como o problema continua em aberto, evidenciando a complexidade do problema.


\newpage

\begin{thebibliography}{99}

\bibitem{putnam1980}
Hilary Putnam,
\textit{Models and Reality},
Journal of Symbolic Logic, Vol. 45, No. 3, 1980, pp. 464--482.

\bibitem{devitt1984}
Michael Devitt,
\textit{Realism and Truth},
Princeton University Press, 1984.

\bibitem{shapiro1997}
Stewart Shapiro,
\textit{Philosophy of Mathematics: Structure and Ontology},
Oxford University Press, 1997.

\bibitem{benacerraf1965}
Benacerraf, P. (1965).
``What Numbers Could Not Be''.
In P. Benacerraf \& H. Putnam (Eds.),
\textit{Philosophy of Mathematics: Selected Readings}.
Cambridge University Press, 1984.

\bibitem{buttonwalsh2018}
Button, T. \& Walsh, S. (2018).
\textit{Philosophy and Model Theory}.
Oxford University Press.
(Part A, Chapters 1--2; Part B, Chapters 6--7).

\bibitem{horsten}
Horsten, L.
``Philosophy of Mathematics''.
In \textit{Stanford Encyclopedia of Philosophy},
Section 5.2.

\bibitem{mcgee2006}
McGee, V.
``There's a Rule for Everything''.
In A. Rayo \& G. Uzquiano (Eds.),
\textit{Absolute Generality}.
Oxford University Press, 2006.

\bibitem{murzi-topey2021}
Murzi, J. \& Topey, B. (2021).
``Categoricity by Convention''.
\textit{Philosophical Studies}.

\end{thebibliography}
\end{document}
